\subsection{Caso de uso: Registrar Lote de Semillas}

\textbf{Descripción:} El usuario registra un nuevo lote de semillas con su información completa, asignándole los tipos de análisis que se realizarán.

\vspace{0.5em}
\noindent\textbf{PreCondiciones:}
\begin{itemize}[nosep]
    \item El usuario debe estar autenticado como Analista o Administrador.
    \item Debe existir al menos un cultivar en el sistema.
    \item Debe existir al menos una empresa registrada.
    \item Debe existir al menos un cliente registrado.
\end{itemize}

\vspace{0.5em}
\noindent\textbf{PostCondiciones:}
\begin{itemize}[nosep]
    \item Lote creado con estado activo = true.
    \item Lote disponible para crear análisis de los tipos asignados.
    \item Lote visible en el listado de lotes activos.
\end{itemize}

\vspace{0.5em}
\noindent\textbf{Actores:} Analista, Administrador.

\vspace{0.5em}
\noindent\textbf{Flujo de eventos principal:}
\begin{enumerate}[nosep]
    \item El usuario accede a la página ``Registro de Lotes''.
    \item Completa los datos obligatorios:
    \begin{itemize}[nosep]
        \item Ficha (campo único)
        \item Nombre del lote (campo único)
        \item Cultivar (selección desde catálogo)
        \item Tipo de lote (INTERNO/EXTERNO)
        \item Empresa (selección desde lista de contactos)
        \item Cliente (selección desde lista de contactos)
        \item Fecha de recibo
        \item Kilos limpios
    \end{itemize}
    \item Completa datos opcionales:
    \begin{itemize}[nosep]
        \item Código CC, Código FF
        \item Fecha de entrega
        \item Depósito, Unidad de embolsado, Remitente
        \item Datos de humedad (tipo y valor)
        \item Número de artículo, Origen, Estado
        \item Fecha de cosecha
        \item Observaciones
    \end{itemize}
    \item Selecciona los tipos de análisis que se realizarán al lote:
    \begin{itemize}[nosep]
        \item Pureza
        \item PMS (Peso de Mil Semillas)
        \item Germinación
        \item Tetrazolio
        \item DOSN (Determinación de Otras Semillas en Número)
    \end{itemize}
    \item El sistema valida que:
    \begin{itemize}[nosep]
        \item La ficha no esté duplicada
        \item El nombre del lote no esté duplicado
        \item La fecha de recibo no sea futura
        \item La fecha de cosecha no sea futura
    \end{itemize}
    \item El sistema crea el lote con \texttt{activo = true}.
    \item El lote aparece en la lista de ``Últimos lotes registrados''.
    \item El sistema muestra mensaje de confirmación.
\end{enumerate}

\vspace{0.5em}
\noindent\textbf{Flujos alternativos:}

\vspace{0.3em}
\noindent\textit{1. Ficha duplicada}
\begin{itemize}[nosep]
    \item El sistema valida en tiempo real
    \item Muestra alerta: ``Esta ficha ya está registrada. Por favor, utiliza una ficha diferente''
    \item No permite guardar hasta corregir
\end{itemize}

\vspace{0.3em}
\noindent\textit{2. Nombre de lote duplicado}
\begin{itemize}[nosep]
    \item El sistema valida en tiempo real
    \item Muestra alerta: ``Este nombre de lote ya está registrado. Por favor, utiliza un nombre diferente''
    \item No permite guardar hasta corregir
\end{itemize}

\vspace{0.3em}
\noindent\textit{3. Fecha futura}
\begin{itemize}[nosep]
    \item El sistema valida en backend
    \item Muestra error: ``La fecha de recibo/cosecha no puede ser posterior a la fecha actual''
    \item No permite guardar hasta corregir
\end{itemize}

\vspace{0.3em}
\noindent\textit{4. Sin tipos de análisis asignados}
\begin{itemize}[nosep]
    \item El lote se crea sin restricciones
    \item No se podrán crear análisis hasta editar el lote y asignar tipos
\end{itemize}

\vspace{0.5em}
\noindent\textbf{Requerimientos especiales:}
\begin{itemize}[nosep]
    \item Validación asíncrona de campos únicos (ficha y nombre)
    \item Eliminación automática de duplicados en tipos de análisis asignados
    \item Los tipos de análisis asignados se almacenan como enum JSON
\end{itemize}
